\begin{proof}

\begin{enumerate}[label={$\langle1\rangle$\arabic*.},leftmargin=*]
\item Choose coordinates $(x,y,z)$ with $C_0=(x=y=0)$.
Generic normal multiplicity three gives
$f\in(x,y)^{(3)}$. Powers of this smooth
complete-intersection ideal are primary, so
$f\in(x,y)^3$ at the specified point.

\item Suppose its nonzero cubic normal tensor is a cube.
A linear normal-coordinate change makes it $x^3$, up
to a nonzero scalar. Weierstrass preparation in $x$
gives the exact equality
\[
 f=U(x,y,z)\bigl[x^3+a_2(y,z)x^2+a_1(y,z)x+a_0(y,z)\bigr],
\]
where $U$ is a holomorphic unit. The bracket belongs to
$(x,y)^3$, so $a_2,a_1,a_0$ are divisible by $y,y^2,y^3$,
respectively. Depressing the monic cubic is a holomorphic
coordinate change preserving $(x,y)$ and yields
\[
 f=U(x,y,z)\bigl[x^3+y^2B(y,z)x+y^3C(y,z)\bigr],
 \qquad B(0,0)=C(0,0)=0.                         \tag{R19.1}
\]
The unit remains in this equality and in the scaling
comparison below.

\item Set
\[
 k=\operatorname{ord}_z B(0,z),\qquad
 l=\operatorname{ord}_z C(0,z),\qquad
 \rho=\min(k/2,l/3),
\]
where the order of an identically zero series is infinity.
At least one order is finite: otherwise all generic normal
cubics would be $x^3$, contrary to generic $D_4$.
Thus $0<\rho<\infty$. Choose $L>3\rho$ and positive
rational weights
\[
       \operatorname{wt}(x)=L+\rho,\qquad
       \operatorname{wt}(y)=L,\qquad
       \operatorname{wt}(z)=1.
\]
Multiplication by a common denominator gives positive
integer weights. The weighted degree of the first face
is $3L+3\rho$. Terms of $B$ containing $y$ have weight
gap at least $L-2\rho>0$, and those of $C$ containing
$y$ have gap at least $L-3\rho>0$. Consequently the
leading polynomial, apart from the nonzero factor $U(0)$,
is exactly one of
\[
\begin{array}{ll}
 P=x(x^2+a y^2z^k),\quad a\ne0,
       & k/2<l/3,\\[2pt]
 P=x^3+b y^3z^l,\quad b\ne0,
       & l/3<k/2,\\[2pt]
 P=Q(x,yz^r),\quad
 Q(X,V)=X^3+aXV^2+bV^3,\quad a,b\ne0,
       & k=2r, l=3r, r\ge1.
\end{array}                                                   \tag{R19.2}
\]
In the last case $r$ is an integer. The depressed cubic
$Q$ is not a cube: a cube with $X^3$ coefficient one and
$X^2V$ coefficient zero would have $a=b=0$. A double
factor of $Q$ is allowed throughout the argument.

\par\medskip\noindent\textbf{The compact class used in the comparison.}
Let $B$ be a connected smooth complex curve and let
$h:B\to\mathbb C$ have two distinct simple zeros $p,q$.
Join them by a smooth embedded arc $\gamma$ whose
interior contains no zero. On the smooth surface
\[
             Z_h=\{(b,u,v)\in B\times\mathbb C^2:
                                      uv=h(b)\},
\]
consider the locus over $\gamma$ with $|u|=|v|$.
It is a circle over each interior point and collapses
to a point at each end. Near an endpoint choose the
arc so that $h$ is a fixed nonzero complex multiple
of a positive real parameter $s$. Then
$u=\sqrt{s}\,e^{\mathrm i\phi}$ and
$v=c\sqrt{s}\,e^{-\mathrm i\phi}$ are the smooth
polar coordinates of a real disc in the local $uv$
chart. The locus is thus a smoothly embedded sphere.

It is Lagrangian for a product K\"ahler form.
Writing $u=R e^{\mathrm i\phi}$,
$v=R e^{\mathrm i(\psi-\phi)}$, the sum of the two
fibre area forms restricts to
$d(R^2/2)\wedge d\psi=0$, since $R,\psi$ depend
only on the arc parameter. The base area form also
restricts to zero. Multiplication by $\mathrm i$
identifies tangent and normal bundles, with the
opposite normal orientation relative to the complex
ambient orientation. Its self-intersection is
therefore $-\chi(S^2)=-2$.
For the compactly supported Poincar\'e dual $\alpha$,
\[
                         \int_{Z_h}\alpha^2=-2.             \tag{R19.3}
\]

If $\pi:W\to Z_h$ is a finite proper holomorphic map
of degree $d>0$ between smooth complex surfaces, then
$\pi^*\alpha$ is compactly supported and
\[
              \int_W(\pi^*\alpha)^2
                       =d\int_{Z_h}\alpha^2=-2d.            \tag{R19.4}
\]
The equality follows from the proper oriented degree;
off the analytic branch locus the map has $d$ positive
sheets, and the branch locus has measure zero for
integration of compactly supported top forms.
Poincar\'e duality gives a compact homology two-class
with that square. Its representative need not itself
be a sphere.

\par\medskip\noindent\textbf{The exact finite maps.}
In the first case of (R19.2), $P=1$ forces $x\ne0$
and gives
\[
           y^2z^k=h(x),\qquad
           h(x)=\frac{1-x^3}{ax},\qquad x\in\mathbb C^*.
\]
The map
\[
 (x,y,z)\longmapsto(x,u=y^2,v=z^k)
\]
is finite proper of degree $2k$ onto $uv=h(x)$.
The target is smooth: at $u=v=0$ the zeros of $h$
are simple, and elsewhere a derivative in $u$ or
$v$ is nonzero. The source is smooth by the weighted
Euler identity on $P=1$. Join two of the three cube
roots of one by an arc in $\mathbb C^*$ avoiding the
third. Equations (R19.3)--(R19.4) give a compact
two-class of square $-4k$ on the leading fibre.

In the second case the exact finite proper map is
\[
 \begin{split}
 \{x^3+b y^3z^l=1\}&\longrightarrow
                 \{uv=(1-x^3)/b\},\\
 (x,y,z)&\longmapsto(x,u=y^3,v=z^l).
 \end{split}
\]
It has degree $3l$, and the base is $\mathbb C$.
The same three simple zeros supply a pulled-back
compact class of square $-6l$.

For the balanced case let
\[
                 B_Q=\{(X,V)\in\mathbb C^2:Q(X,V)=1\}.
\]
It is smooth by the ordinary Euler identity and is
connected even when $Q$ has a double factor. Its
projection to $\mathbb P^1$ minus the distinct zero
directions of $Q$ is an unramified cyclic cover of
degree three. Around a zero direction of multiplicity
one or two its monodromy generates $\mathbb Z/3$.
Such a direction exists because $Q$ is not a cube,
so the cover is connected.

The function $V$ on $B_Q$ has precisely the three
zeros $(\zeta,0)$ with $\zeta^3=1$, and they are
simple because $Q_X(\zeta,0)=3\zeta^2\ne0$.
The surface
\[
             Z_Q=\{(b,u,w):b\in B_Q,\ uw=V(b)\}
\]
is smooth. The exact map
\[
 \begin{split}
       \{Q(x,yz^r)=1\}&\longrightarrow Z_Q,\\
       (x,y,z)&\longmapsto
                    \bigl(b=(x,yz^r),u=y,w=z^r\bigr)
 \end{split}                                                   \tag{R19.5}
\]
is finite proper of degree $r$: $z$ satisfies the
monic equation $z^r=w$ and the other coordinates are
retained. An arc in $B_Q$ joining two $V$-zeros and
avoiding the third supplies a compact class of square
$-2r$ on the source. In particular this argument
retains $Q=(X-\lambda V)^2(X+2\lambda V)$,
$\lambda\ne0$, without a perturbation of its
repeated factor.

\par\medskip\noindent\textbf{Transport to the actual local fibre.}
Let $(w_x,w_y,w_z)$ be the positive integer weights
and let $d$ be the resulting degree. Put
\[
 D_\epsilon(x,y,z)=
       (\epsilon^{w_x}x,\epsilon^{w_y}y,\epsilon^{w_z}z),
 \qquad
 G_\epsilon=\epsilon^{-d}f\circ D_\epsilon.
\]
These actual functions converge with derivatives on
compact sets to $U(0)P$. The weighted Euler identity
makes $dP$ nonzero on $P=1$, including the repeated-factor
case on its zero set. Choose compact representatives
of the two-class above and of the transverse
perturbation computing its square. On a compact
neighbourhood of these chains, the gradient-flow
comparison of P4.52 solves
\[
       dG_\epsilon(V_\epsilon)
                         =-\partial_\epsilon G_\epsilon.
\]
The denominator is bounded away from zero on a
slightly larger neighbourhood; a cutoff gives a
flow for sufficiently small positive $\epsilon$.
It transports the chains and their intersections
to $G_\epsilon=U(0)$. Scaling by $D_\epsilon$ then
embeds them in the original fixed Milnor ball on
\[
                         f=U(0)\epsilon^d.
\]
The oriented intersection number is preserved and
equals $-4k$, $-6l$, or $-2r$. Each is nonzero,
contradicting P4.51. This proves the assertion.

\end{enumerate}
\end{proof}
